% Options for packages loaded elsewhere
\PassOptionsToPackage{unicode}{hyperref}
\PassOptionsToPackage{hyphens}{url}
\documentclass[
]{article}
\usepackage{xcolor}
\usepackage[margin=1in]{geometry}
\usepackage{amsmath,amssymb}
\setcounter{secnumdepth}{-\maxdimen} % remove section numbering
\usepackage{iftex}
\ifPDFTeX
  \usepackage[T1]{fontenc}
  \usepackage[utf8]{inputenc}
  \usepackage{textcomp} % provide euro and other symbols
\else % if luatex or xetex
  \usepackage{unicode-math} % this also loads fontspec
  \defaultfontfeatures{Scale=MatchLowercase}
  \defaultfontfeatures[\rmfamily]{Ligatures=TeX,Scale=1}
\fi
\usepackage{lmodern}
\ifPDFTeX\else
  % xetex/luatex font selection
\fi
% Use upquote if available, for straight quotes in verbatim environments
\IfFileExists{upquote.sty}{\usepackage{upquote}}{}
\IfFileExists{microtype.sty}{% use microtype if available
  \usepackage[]{microtype}
  \UseMicrotypeSet[protrusion]{basicmath} % disable protrusion for tt fonts
}{}
\makeatletter
\@ifundefined{KOMAClassName}{% if non-KOMA class
  \IfFileExists{parskip.sty}{%
    \usepackage{parskip}
  }{% else
    \setlength{\parindent}{0pt}
    \setlength{\parskip}{6pt plus 2pt minus 1pt}}
}{% if KOMA class
  \KOMAoptions{parskip=half}}
\makeatother
\usepackage{longtable,booktabs,array}
\usepackage{calc} % for calculating minipage widths
% Correct order of tables after \paragraph or \subparagraph
\usepackage{etoolbox}
\makeatletter
\patchcmd\longtable{\par}{\if@noskipsec\mbox{}\fi\par}{}{}
\makeatother
% Allow footnotes in longtable head/foot
\IfFileExists{footnotehyper.sty}{\usepackage{footnotehyper}}{\usepackage{footnote}}
\makesavenoteenv{longtable}
\usepackage{graphicx}
\makeatletter
\newsavebox\pandoc@box
\newcommand*\pandocbounded[1]{% scales image to fit in text height/width
  \sbox\pandoc@box{#1}%
  \Gscale@div\@tempa{\textheight}{\dimexpr\ht\pandoc@box+\dp\pandoc@box\relax}%
  \Gscale@div\@tempb{\linewidth}{\wd\pandoc@box}%
  \ifdim\@tempb\p@<\@tempa\p@\let\@tempa\@tempb\fi% select the smaller of both
  \ifdim\@tempa\p@<\p@\scalebox{\@tempa}{\usebox\pandoc@box}%
  \else\usebox{\pandoc@box}%
  \fi%
}
% Set default figure placement to htbp
\def\fps@figure{htbp}
\makeatother
\setlength{\emergencystretch}{3em} % prevent overfull lines
\providecommand{\tightlist}{%
  \setlength{\itemsep}{0pt}\setlength{\parskip}{0pt}}
\usepackage{booktabs}
\usepackage{longtable}
\usepackage{array}
\usepackage{multirow}
\usepackage{wrapfig}
\usepackage{float}
\usepackage{colortbl}
\usepackage{pdflscape}
\usepackage{tabu}
\usepackage{threeparttable}
\usepackage{threeparttablex}
\usepackage[normalem]{ulem}
\usepackage{makecell}
\usepackage{xcolor}
\usepackage{multicol}
\usepackage{ulem}
\usepackage{hhline}
\newlength\Oldarrayrulewidth
\newlength\Oldtabcolsep
\usepackage{hyperref}
\usepackage{bookmark}
\IfFileExists{xurl.sty}{\usepackage{xurl}}{} % add URL line breaks if available
\urlstyle{same}
\hypersetup{
  pdftitle={Withdraw models and notes},
  pdfauthor={Bill Prisbrey},
  hidelinks,
  pdfcreator={LaTeX via pandoc}}

\title{Withdraw models and notes}
\author{Bill Prisbrey}
\date{2025-12-04}

\begin{document}
\maketitle

\textbf{PURPOSE:} \emph{The purpose of this report is to explore
`Withdraw' grades using different predictive models.}

\section{Executive Summary}\label{executive-summary}

This report finds that ``Withdraw'' grades are a rare phenomenon with
essentially no relationship to pre-admit data.

\global\setlength{\Oldarrayrulewidth}{\arrayrulewidth}

\global\setlength{\Oldtabcolsep}{\tabcolsep}

\setlength{\tabcolsep}{2pt}

\renewcommand*{\arraystretch}{1.5}



\providecommand{\ascline}[3]{\noalign{\global\arrayrulewidth #1}\arrayrulecolor[HTML]{#2}\cline{#3}}

\begin{longtable}[c]{|p{0.60in}|p{1.00in}|p{0.80in}|p{4.00in}}

\caption{Predictive\ model\ accuracy\ is\ essentially\ random.}\\

\ascline{1.5pt}{666666}{1-4}

\multicolumn{1}{>{\raggedright}m{\dimexpr 0.6in+0\tabcolsep}}{\textcolor[HTML]{000000}{\fontsize{11}{11}\selectfont{ver}}} & \multicolumn{1}{>{\raggedright}m{\dimexpr 1in+0\tabcolsep}}{\textcolor[HTML]{000000}{\fontsize{11}{11}\selectfont{Description}}} & \multicolumn{1}{>{\raggedleft}m{\dimexpr 0.8in+0\tabcolsep}}{\textcolor[HTML]{000000}{\fontsize{11}{11}\selectfont{Kappa}}} & \multicolumn{1}{>{\raggedright}m{\dimexpr 4in+0\tabcolsep}}{\textcolor[HTML]{000000}{\fontsize{11}{11}\selectfont{Notes}}} \\

\ascline{1.5pt}{666666}{1-4}\endfirsthead \caption[]{Predictive\ model\ accuracy\ is\ essentially\ random.}\\

\ascline{1.5pt}{666666}{1-4}

\multicolumn{1}{>{\raggedright}m{\dimexpr 0.6in+0\tabcolsep}}{\textcolor[HTML]{000000}{\fontsize{11}{11}\selectfont{ver}}} & \multicolumn{1}{>{\raggedright}m{\dimexpr 1in+0\tabcolsep}}{\textcolor[HTML]{000000}{\fontsize{11}{11}\selectfont{Description}}} & \multicolumn{1}{>{\raggedleft}m{\dimexpr 0.8in+0\tabcolsep}}{\textcolor[HTML]{000000}{\fontsize{11}{11}\selectfont{Kappa}}} & \multicolumn{1}{>{\raggedright}m{\dimexpr 4in+0\tabcolsep}}{\textcolor[HTML]{000000}{\fontsize{11}{11}\selectfont{Notes}}} \\

\ascline{1.5pt}{666666}{1-4}\endhead



\multicolumn{1}{>{\raggedright}m{\dimexpr 0.6in+0\tabcolsep}}{\textcolor[HTML]{000000}{\fontsize{11}{11}\selectfont{vA0}}} & \multicolumn{1}{>{\raggedright}m{\dimexpr 1in+0\tabcolsep}}{\textcolor[HTML]{000000}{\fontsize{11}{11}\selectfont{Admissions\ data}}} & \multicolumn{1}{>{\raggedleft}m{\dimexpr 0.8in+0\tabcolsep}}{\textcolor[HTML]{000000}{\fontsize{11}{11}\selectfont{0.025}}} & \multicolumn{1}{>{\raggedright}m{\dimexpr 4in+0\tabcolsep}}{\textcolor[HTML]{000000}{\fontsize{11}{11}\selectfont{'propensity\ to\ withdraw'\ is\ not\ found\ in\ admit\ data}}} \\





\multicolumn{1}{>{\raggedright}m{\dimexpr 0.6in+0\tabcolsep}}{\textcolor[HTML]{000000}{\fontsize{11}{11}\selectfont{vB0}}} & \multicolumn{1}{>{\raggedright}m{\dimexpr 1in+0\tabcolsep}}{\textcolor[HTML]{000000}{\fontsize{11}{11}\selectfont{Adds\ credit\ load}}} & \multicolumn{1}{>{\raggedleft}m{\dimexpr 0.8in+0\tabcolsep}}{\textcolor[HTML]{000000}{\fontsize{11}{11}\selectfont{0.037}}} & \multicolumn{1}{>{\raggedright}m{\dimexpr 4in+0\tabcolsep}}{\textcolor[HTML]{000000}{\fontsize{11}{11}\selectfont{suggests\ concurrent\ information\ is\ slightly\ more\ important}}} \\





\multicolumn{1}{>{\raggedright}m{\dimexpr 0.6in+0\tabcolsep}}{\textcolor[HTML]{000000}{\fontsize{11}{11}\selectfont{vC0}}} & \multicolumn{1}{>{\raggedright}m{\dimexpr 1in+0\tabcolsep}}{\textcolor[HTML]{000000}{\fontsize{11}{11}\selectfont{Adds\ GPA\ prediction}}} & \multicolumn{1}{>{\raggedleft}m{\dimexpr 0.8in+0\tabcolsep}}{\textcolor[HTML]{000000}{\fontsize{11}{11}\selectfont{0.029}}} & \multicolumn{1}{>{\raggedright}m{\dimexpr 4in+0\tabcolsep}}{\textcolor[HTML]{000000}{\fontsize{11}{11}\selectfont{'pred'\ is\ new\ top\ variable;\ suggests\ data\ leak}}} \\





\multicolumn{1}{>{\raggedright}m{\dimexpr 0.6in+0\tabcolsep}}{\textcolor[HTML]{000000}{\fontsize{11}{11}\selectfont{vD0}}} & \multicolumn{1}{>{\raggedright}m{\dimexpr 1in+0\tabcolsep}}{\textcolor[HTML]{000000}{\fontsize{11}{11}\selectfont{Rolling\ GPA\ prediction}}} & \multicolumn{1}{>{\raggedleft}m{\dimexpr 0.8in+0\tabcolsep}}{\textcolor[HTML]{000000}{\fontsize{11}{11}\selectfont{0.016}}} & \multicolumn{1}{>{\raggedright}m{\dimexpr 4in+0\tabcolsep}}{\textcolor[HTML]{000000}{\fontsize{11}{11}\selectfont{rolling\ prediction\ eliminates\ leak;\ still\ no\ accuracy}}} \\

\ascline{1.5pt}{666666}{1-4}



\end{longtable}



\arrayrulecolor[HTML]{000000}

\global\setlength{\arrayrulewidth}{\Oldarrayrulewidth}

\global\setlength{\tabcolsep}{\Oldtabcolsep}

\renewcommand*{\arraystretch}{1}

\section{Discussion}\label{discussion}

Kappa values of \textless0.05 suggest that the propensity to withdraw is
simply not found in data available at the time of admission. In other
words, the ACT test doesn't have a `propensity to withdraw' section.

Adding the concurrent credit load improved the Kappa. It's so low it
might just be random noise. Or, it might be a suggestion that
concurrent, contemporaneous, or situational information might be more
predictive of a withdraw than admit-level data.

Adding the GPA prediction resulted in that prediction becoming the new
top variable, but the accuracy dropped. Again, the values are so low it
could just be noise. Or, it could be a data leak in that I applied the
prediction to the same data that I trained on.

Substituting a rolling GPA prediction to eliminate a data leak resulted
in an even lower Kappa. (However, a closer investigation revealed a
discontinuity between prior years and year 2024. This was significant at
the p\textless.1 level per a logistic regression model. This is a
whisper that there might be a detectable relationship between `W' and
predicted grades going forward, or it might be more noise.)

The fact that the admissions data was able to predict grades with some
accuracy, but couldn't distinguish withdrawals, suggests that withdraws
operate by an entirely different mechanism. The addition of predicted
grades without improving the model adds evidence to this conclusion.

It is worthy of its own investigation as a separate project to
brainstorm drivers and seek new data sets.

Potential causes and underlying factors of withdrawals may include
things like:

\begin{itemize}
\tightlist
\item
  \emph{Behavioral tendencies,} where some students have personality
  traits that might more accurately predict a `W'.\\
\item
  \emph{Difficulty adjusting} to university-level studies.\\
\item
  \emph{Personal preference} for an instructional style or a personality
  match between instructor and student.\\
\item
  \emph{Completely external} to the academic experience, such as a
  medical, family, or financial difficulty.\\
\item
  \emph{Schedule difficulties} unrelated to academic abilities may
  precipitate a withdrawal.
\end{itemize}

New data sets could include:

\begin{itemize}
\tightlist
\item
  \emph{Early grades} for homework, quizzes, and tests.\\
\item
  \emph{Schedule} to determine potential conflicts.
\end{itemize}

With the `admit-level' information used in these models, and the
predicted grade in vD0, it does not seem that `withdraw' students are
adjacent to flunking students, or related to the expected grade at all,
otherwise the models would be more accurate.

A few avenues I would like to investigate includes:

\begin{itemize}
\tightlist
\item
  \emph{Match:} Is a student more likely to withdraw from a math course
  not appropriate to their abilities?

  \begin{itemize}
  \tightlist
  \item
    Use the model to predict an expected grade, and incorporate that as
    a prediction variable. Lower expected grades should have higher
    withdraw rates.\\
  \item
    Develop counterfactuals and rank the courses. The lower rankings
    (lower expected GPA) should have higher withdraw rates.\\
  \item
    Develop a `difficulty' ranking of the courses and use as a
    predictive variable.
  \item
    Use the catalog number itself as a difficulty ranking\\
  \end{itemize}
\item
  \textbf{\emph{Result of investigating match (vD0):}} Version D0 used a
  rolling prediction of grade based on prior years to avoid a data leak.
  Kappa was still very low. This suggests that withdrawals are unrelated
  to academic ability and that the alternative approaches described
  above should be pursued.

  \begin{itemize}
  \tightlist
  \item
    However, it is worth noting that the year 2024 saw an improvement in
    the relationship to predicted grades to a significance of p
    \textless{} .1. An investigation of the distribution saw that higher
    grades were actually more likely to withdraw. A possible explanation
    is that in 2024, students who submitted ACT math scores are more
    elite and therefore more cautious of their grades. If these students
    were anticipating a ``B'' instead of an ``A'' then they might
    withdraw. Or in other words, the behavior of the subset of students
    used in this model for 2024 actually do have a relationship to
    predicted grades and are different than the behavior of the general
    population from years prior.
  \end{itemize}
\end{itemize}

Because vD0 failed to improve Kappa, a separate project is suggested.

\section{Next steps}\label{next-steps}

\begin{itemize}
\tightlist
\item
  Abandon using the current approach of admit-level data and pursue a
  separate project as described above.
\end{itemize}

Possible steps if continuing the current approach of using admit- and
course- level data:

\begin{itemize}
\item
  Create a rolling prediction of GPA values from models trained on the
  prior data. \textbf{\emph{DONE}}
\item
  Rank the counterfactual GPA's per student and use as a variable.
\item
  Approximate course difficulty

  \begin{itemize}
  \tightlist
  \item
    With the catalog number
  \item
    With a novel approach (by developing the `gentle guidelines'
    suggested by the predictive models)
  \end{itemize}
\end{itemize}

\section{APPENDIX}\label{appendix}

\subsubsection{Withdraw is rare or highly
imbalanced}\label{withdraw-is-rare-or-highly-imbalanced}

\pandocbounded{\includegraphics[keepaspectratio]{Withdraw-models-and-notes_files/figure-latex/unnamed-chunk-6-1.pdf}}

\subsubsection{Withdraw does not appear related to expected
grades}\label{withdraw-does-not-appear-related-to-expected-grades}

\pandocbounded{\includegraphics[keepaspectratio]{Withdraw-models-and-notes_files/figure-latex/unnamed-chunk-7-1.pdf}}

\subsubsection{Withdraw may be slightly related to expected
grades}\label{withdraw-may-be-slightly-related-to-expected-grades}

\pandocbounded{\includegraphics[keepaspectratio]{Withdraw-models-and-notes_files/figure-latex/unnamed-chunk-8-1.pdf}}

\subsubsection{vA0 predictors}\label{va0-predictors}

\begin{longtable}[]{@{}l@{}}
\toprule\noalign{}
x \\
\midrule\noalign{}
\endhead
\bottomrule\noalign{}
\endlastfoot
course \\
SEX \\
FIRST\_GEN\_STATUS\_CD \\
ETHNICITY \\
RESSTAT \\
FA\_PELL \\
APCREDIT \\
HSGPA \\
HSPRIVATE \\
HONORS \\
ACTCOMP \\
ACTENGL \\
ACTMATH \\
ACTSCI \\
cohort\_year \\
class\_year \\
yr\_diff \\
season \\
course\_level \\
vol\_cluster \\
age\_cut \\
\end{longtable}

\subsubsection{vB0 predictors}\label{vb0-predictors}

Add ``load'' as the number of credits taken

\subsubsection{vC0 predictors}\label{vc0-predictors}

Add ``pred'' as the predicted GPA from a model trained on years less
than or equal to 2023, and applied to all years.

\subsubsection{vD0 predictors}\label{vd0-predictors}

Substitute ``pred'' as the predicted GPA from rolling models trained on
years prior to the year predicted.

\end{document}
